%laden der Präambel mit Latexbefehlen/-klassen
\include{expose.preambel}

\subject{Exposé zur Abschlussarbeit}
\title{Aussagekräftiger Titel der Abschlussarbeit}
\author[1,2]{Jane Doe}
\author[2]{externer Betreuer}
\author[1]{Mathias Hofmann}
\affil[1]{TU Berlin, Institut für Energietechnik, Marchstraße 18, 10587 Berlin}
\affil[2]{Firma, Adresse}
\date{\today}

\begin{document}

\maketitle

\section*{Exposé}

Ein Exposé soll die folgenden sieben Kernpunkte beschreiben. Deren Gewichtung kann sich je nach Aufgabenstellung unterscheiden. Bei Unklarheiten liefern Ebel und Bliefert~\cite{ebel2009} eine für das wissenschaftliche Arbeiten ergiebige Quelle.

\begin{itemize}
    \item Problembereich
    \item Forschungsstand
    \item Fragestellungen
    \item Methoden
    \item Materialzugang
    \item Arbeitsplan
    \item Ausgewählte Literatur
\end{itemize}

\section*{Gliederung}

Erstellen Sie eine Gliederung, um den Umfang möglicher Arbeitspakete besser abschätzen zu können. Insbesondere im Hinblick auf den schriftlichen Teil ihrer Arbeit.

\begin{enumerate}
    \item Einleitung
    \item Grundlagen, Stand der Forschung, Recherche
    \item Methodik
        \begin{enumerate}
            \item Eine tolle Methode
            \item Noch eine super Methode
        \end{enumerate}
    \item Analyse, Berechnung, Modell, Annahmen
    \item Auswertung, Ergebnisse, Diskussion
    \item Fazit und Ausblick
\end{enumerate}

\section*{Zeitplan}

Mit dem Paket pgfgantt können Gantt-Charts direkt im Latex-Quellcode erstellt und eingebunden werden. Diese eignen sich hervorragend zur Visualisierung eines Zeitplans.

\begin{center}
    \input{bilder/gantt-chart}    
\end{center}


%Literatur
\printbibliography

\end{document}
